% Generated draft: requires mathematical and source review.
\documentclass[11pt]{article}
\usepackage[margin=25mm]{geometry}
\usepackage{fontspec}
\setmainfont{Latin Modern Roman}
\usepackage{amsmath,amssymb,mathtools,unicode-math}
\setmathfont{Latin Modern Math}
\usepackage{enumitem,needspace,xurl,hyperref}
\setlength{\parindent}{0pt}
\setlength{\parskip}{5pt}
\setlength{\emergencystretch}{3em}
\newcommand{\N}{\mathbb{N}}
\newcommand{\Z}{\mathbb{Z}}
\newcommand{\Q}{\mathbb{Q}}
\newcommand{\R}{\mathbb{R}}
\newcommand{\C}{\mathbb{C}}
\newcommand{\F}{\mathbb{F}}
\newcommand{\A}{\mathbb{A}}
\newcommand{\PP}{\mathbb{P}}
\newcommand{\E}{\mathbb{E}}
\newcommand{\eps}{\varepsilon}
\newcommand{\dd}{\,\mathrm{d}}
\begin{document}
\section*{UnsolvedMath ID: 30006443}
\textbf{Source problem code:} catalogue entry
\subsection{Definitions and mathematical statement}
Write \(\tau(V,0)\) for the Zariski dimensionality type of a singularity germ \((V,0)\). This notation names the invariant referred to in the record; its definition and the class of germs on which it is defined are not supplied.

The damaged source excerpt indicates projections of the schematic form
\[
\pi\colon (k^{r+1},0)\longrightarrow (k^r,0).
\]
For an admissible projection, write \(\Delta_\pi\) for the discriminant locus of its restriction to the germ. Here ``admissible'' and ``discriminant'' refer to the intended Zariski projection construction; the record does not specify their hypotheses or conventions.

\textbf{Question (with unresolved source qualifications).} For every singularity germ \((V,0)\) in the intended class satisfying
\[
\tau(V,0)\geq 3,
\]
is a generic linear projection generic in the sense required by Zariski's construction, so that its discriminant \(\Delta_\pi\) can be used to compute \(\tau(V,0)\)?

The scope is all dimensionality types at least three. The record does not supply the definition of dimensionality type, its discriminant computation rule, the ground-field hypotheses, the precise category of singularities, or the parameter space and exceptional set defining a generic linear projection. Consequently, this is a qualified statement of the intended question, not a fully specified quantified conjecture. In particular, no numerical recurrence relating the dimensionality types of \(V\) and \(\Delta_\pi\) is asserted here.
\subsection{Short English statement}
The question asks whether generic linear projections suffice for the discriminant construction of Zariski dimensionality type for every singularity of type at least three. The supplied record labels this question open, but essential definitions and hypotheses are missing, so the formulation requires source consultation and human mathematical review.
\subsection{Sources}
\begin{itemize}
\item Local UnsolvedMath registry, permanent ID \(30006443\), title \emph{Generic Linear Projections and Dimensionality Type}, source problem code \texttt{catalogue entry}; supplied complete source record. No individual proposer is identified.
\item Supplied report citation: \emph{Singularities (2026)}, Oberwolfach Reports, DOI \texttt{10.4171/owr/2025/45}, pp.\ 2413--2490. \url{https://doi.org/10.4171/owr/2025/45}. The same report is listed among the supplied literature sources as \emph{Oberwolfach Report 45/2025}.
\item Supplied literature source: \emph{Zariski dimensionality type}. \url{https://math.univ-cotedazur.fr/~parus/publis/JAG1439R1.pdf}. Authors and further bibliographic details are not supplied.
\item Literature assessment dated exactly \texttt{2026-08-22}: the supplied curation labels the problem open and reports that the generic-polynomial projection theorem and the dimensionality-type-two case are known, while sufficiency of generic linear projection in type three or higher remains open. This is supplied dated curation, not an independently verified literature assessment; the record requests independent verification of time-sensitive status and cited claims.
\item The registry labels the statement \texttt{corrected\_verified} and describes its changes as repairs to formula delimiters, notation, grammar, and typography that preserve mathematical content and scope. These are supplied editorial labels, not independent mathematical verification.
\end{itemize}
\end{document}
